\documentclass[11pt]{article}
\usepackage{amsmath,amsthm,amssymb}
\usepackage[margin=1in]{geometry}
\usepackage{booktabs}
\usepackage{array}
\usepackage{stmaryrd}
\usepackage[colorlinks=true,linkcolor=blue,citecolor=blue]{hyperref}

\newtheorem{theorem}{Theorem}
\newtheorem{proposition}[theorem]{Proposition}
\newtheorem{corollary}[theorem]{Corollary}
\newtheorem{lemma}[theorem]{Lemma}
\theoremstyle{definition}
\newtheorem{definition}[theorem]{Definition}
\newtheorem{remark}[theorem]{Remark}
\newtheorem{example}[theorem]{Example}
\newtheorem{problem}[theorem]{Problem}

\newcommand{\SPoly}{\mathsf{SPoly}}
\newcommand{\PolyN}{\mathsf{Poly}_{\mathbb{N}}}
\newcommand{\Poly}{\mathsf{Poly}}
\newcommand{\Cont}{\mathsf{Cont}}
\newcommand{\DCont}{\mathsf{DCont}}
\newcommand{\Cat}{\mathsf{Cat}}
\newcommand{\Set}{\mathsf{Set}}
\newcommand{\Fam}{\mathsf{Fam}}
\newcommand{\Mod}{\mathsf{Mod}}
\newcommand{\N}{\mathbb{N}}
\newcommand{\DD}{\mathbb{D}}
\newcommand{\sem}[1]{\llbracket #1 \rrbracket}
\newcommand{\id}{\mathrm{id}}
\newcommand{\rank}{\operatorname{rank}}
\newcommand{\Ob}{\operatorname{Ob}}
\newcommand{\Out}{\operatorname{Out}}
\renewcommand{\ker}{\operatorname{ker}}

\title{Toward a uniqueness theorem for the container derivative\\
\large the solid infinitesimal parts of polynomial functors}
\author{MacBeth\\ \small Kodamai}
\date{3 October 2026}

\begin{document}
\maketitle

\begin{abstract}
A companion note showed that the Abbott--Altenkirch--Ghani--McBride one-hole derivative
$\partial$ is the differential combinator of a Cartesian differential category of polynomial
functors, realised by the dual-number substitution $T(p)=p(y+\varepsilon v)\bmod\varepsilon^2$.
That settles \emph{existence}. This note asks for \emph{uniqueness}: is $\partial$ the only
nontrivial first-order tangent structure a polynomial functor can carry? We prove the decisive
rigidity input. An infinitesimal object $D=y\oplus M$ in the substitution monoidal category
$(\Poly,\lhd,y)$ has, for a tangent bundle $T=(-)\lhd D$, a \emph{solid} infinitesimal part when
its multiplication $M\lhd M\to M$ is invertible; we show that over the rig $\N$ the linear
containers are automatically free, so solidity reads $|S|^2=|S|$, whose only finite solutions are
$0$ and $1$. Hence there are \emph{exactly two} finite-support solid infinitesimal parts, $M=0$
and $M=y$, corresponding to the trivial structure and to $\partial$. The finiteness is sharp:
$M=\N\cdot y$ is a third solid part, since $|\N\times\N|=|\N|$. This is the polynomial-functor
analogue of Lanfranchi--Lemay's rigidity for affine schemes, and cleaner --- no principal-ideal
hypothesis is needed. Uniqueness of the tangent structure itself follows once one knows that every
representable first-order tangent structure has a solid infinitesimal part; we verify this
converse on the two primitives and against the Whitney sums, and state the general case as the one
open problem that stands between this classification and a full uniqueness theorem.
\end{abstract}

\section{Introduction}\label{sec:intro}

\paragraph{From existence to uniqueness.}
In a companion note \cite{MacBethTangent} we showed that the category $\SPoly$ of finitary
polynomial functors is a Cartesian differential category whose differential combinator is the
Abbott--Altenkirch--Ghani--McBride one-hole-context derivative $\partial$ \cite{AAGM2003}, and
that the attached tangent bundle is computed by the oldest trick in differential algebra:
substitute a dual number and read off the $\varepsilon$-coefficient,
\[
  T(p)\ =\ p(y+\varepsilon v)\bmod\varepsilon^2\ =\ p(y)+\varepsilon\,v\cdot\partial p.
\]
That note answers an \emph{existence} question --- there is a tangent structure on polynomial
functors whose vertical part is $\partial$. The natural next question is \emph{uniqueness}. The
dual-number structure is one tangent structure; the identity functor $T=\id$ is another (the
trivial one). Are there others? Could a polynomial functor carry an exotic first-order tangent
bundle that is not the derivative and not trivial?

\paragraph{The template.}
For commutative algebra this question has a sharp answer. Lanfranchi and Lemay
\cite{LanfranchiLemay2025} classify the \emph{representable} tangent structures on the opposite of
the category of modules (equivalently, affine schemes): a representable tangent structure is a
tensoring $T=D\otimes(-)$ by an \emph{infinitesimal object} $D=R\oplus M$, and when the base ring
$R$ is a principal ideal domain there are \emph{exactly two} --- the trivial one ($M=0$) and the
Kähler/dual-number one ($M=R$). The mechanism is a rank computation: the infinitesimal part $M$ of
a tangent structure must be \emph{solid}, $M\otimes_R M\cong M$; over a PID a solid finitely
generated module is free of rank $n$ with $R^{n^2}\cong R^n$, forcing $n^2=n$, hence $n\in\{0,1\}$
(their Theorem~4.17, Corollary~4.18).

\paragraph{The port, and what is genuinely proved.}
This note ports that mechanism to polynomial functors, and the port is in one respect
\emph{cleaner} than the original. The substitution monoidal category $(\Poly,\lhd,y)$ plays the
role of $(\Mod_R,\otimes_R,R)$; an infinitesimal object is a polynomial $D=y\oplus M$ with an
augmentation, and the representable tangent bundle is $T=(-)\lhd D$. The infinitesimal part $M$ is
forced to be a \emph{linear} container $M=S\cdot y$ by the additive-bundle axiom, and here is the
simplification: over the rig $\N$ a linear container \emph{is} a free module, with no
principal-ideal hypothesis to invoke, because the linear polynomial functors are by definition the
free $\N$-modules. Solidity then reads $|S|^2=|S|$ in the cardinal $|S|=\rank M$, whose only
finite solutions are $0$ and $1$. We prove:

\begin{theorem}[Classification of finite solid infinitesimal parts; \S\ref{sec:rigidity}]
\label{thm:main-intro}
A finite-support solid infinitesimal part of $(\Poly,\lhd,y)$ has rank in $\{0,1\}$. There are
exactly two: $M=0$, giving the trivial structure $T=\id$, and $M=y$, giving the AAGM dual-number
structure $T(p)=p(y+\varepsilon v)\bmod\varepsilon^2$ with combinator $\partial$.
\end{theorem}

This is the airtight core, and its arithmetic is machine-checked (\S\ref{sec:rigidity}, a
Mathlib-free Lean~4 oracle with no classical axioms).

\paragraph{The honest gap.}
Theorem~\ref{thm:main-intro} classifies the solid \emph{infinitesimal parts}. To conclude that
there are exactly two tangent \emph{structures} one needs the correspondence
\begin{quote}
$(\star)$\quad representable first-order tangent structure $T=(-)\lhd D$ $\iff$ solid
infinitesimal part $M=\ker(D\to y)$.
\end{quote}
The forward direction (a solid part yields a tangent structure) we verify on the two primitives,
which is all Theorem~\ref{thm:main-intro} leaves to check. The converse (a tangent structure forces
its infinitesimal part solid) is the content of Lanfranchi--Lemay's vertical-lift universality
(their Theorem~4.8), proved over a commutative ring; since $\N$ is a rig, their argument does
\emph{not} transfer verbatim, and we do not reprove it here. We therefore state the resulting
uniqueness as a theorem \emph{conditional on} $(\star_c)$ (the converse), and flag $(\star_c)$ as
the one open problem (\S\ref{sec:correspondence}, Problem~\ref{prob:converse}). We support it with
two pieces of evidence: the forward correspondence on primitives, and the observation that the only
other finite candidates --- the higher square-zero objects $D=y+\sum_{i=1}^r\varepsilon_i v_i$,
$r\ge2$ --- are the Whitney sums $T_r=T\times_X\cdots\times_X T$, which are auxiliary bundles in the
tangent axioms, not standalone structures, and which are indeed non-solid ($\rank=r^2\ne r$). So
solidity coincides with being a primitive tangent structure, exactly as the correspondence
predicts.

\paragraph{The finiteness is sharp.}
Unqualified uniqueness over all of $\Poly$ is false, and we exhibit the obstruction explicitly:
$M=\N\cdot y$ is solid, because $M\lhd M=(\N\times\N)\cdot y\cong\N\cdot y$ by the Cantor pairing
(\S\ref{sec:sharp}, Proposition~\ref{prop:infinite}). This is the precise analogue of
Lanfranchi--Lemay's extra structures when finite generation is dropped (their Examples~4.20--4.21),
and it tells us the right hypothesis is a \emph{finite-dimensional} infinitesimal part --- finitely
many tangent directions --- here taken as plain finiteness of the shape set.

\paragraph{Contribution.}
We do not claim a full uniqueness theorem; we claim its rigidity half. The contributions are: the
classification of finite solid infinitesimal parts (Theorem~\ref{thm:main-intro}), proved over a
rig with no PID hypothesis and Lean-certified in its arithmetic; the single technical observation
that makes the port faithful --- that on linear containers the \emph{bundle} tensor $\lhd$ and the
\emph{module} (Dirichlet) tensor coincide (Lemma~\ref{lem:coincide}), so solidity is computed in the
bundle's own tensor; the sharpness of the finiteness hypothesis
(Proposition~\ref{prop:infinite}); and the precise reduction of full uniqueness to the one ported
converse $(\star_c)$. The result is timely for the tangent-category and \texttt{Poly} communities:
``we have classified the solid infinitesimal parts of finite-support polynomial functors, and
uniqueness of the container-derivative tangent structure reduces to vertical-lift universality'' is
a clean statement of where the problem stands.

\paragraph{Outline.}
\S\ref{sec:prelim} recalls containers, the substitution monoidal structure, infinitesimal objects
and the representable tangent bundle. \S\ref{sec:coincide} proves the bundle/module tensor
coincidence. \S\ref{sec:rigidity} proves the rank rigidity (Theorem~\ref{thm:main-intro}) and
records the Lean certification. \S\ref{sec:sharp} proves the infinite-support obstruction.
\S\ref{sec:correspondence} discusses the correspondence $(\star)$, states the conditional uniqueness
theorem and the open converse. \S\ref{sec:compare} compares with Lanfranchi--Lemay and states the
significance. \S\ref{sec:conclusion} concludes.

\section{Preliminaries}\label{sec:prelim}

We recall only what the rigidity argument uses; for the Cartesian differential and tangent
structure on $\SPoly$ we refer to the companion note \cite{MacBethTangent}.

\paragraph{Containers and the substitution monoidal category.}
A \emph{container in $n$ variables} is a set $S$ of \emph{shapes} together with, for each $s\in S$
and each $i\in\{1,\dots,n\}$, a position set $P^i_s$; its extension is the polynomial functor
$\sem{S\lhd P}(\vec X)=\sum_{s\in S}\prod_{i=1}^n X_i^{P^i_s}$. We work with finitary containers, so
only the cardinalities $|P^i_s|$ matter up to isomorphism \cite{AAGM2003}. Composition of containers
is substitution of extensions, written $G\lhd F$, with unit the identity container $y$; this makes
$(\Poly,\lhd,y)$ a (strict, once we pass to isomorphism classes) monoidal category, and the
comonoids in it are exactly the small categories \cite{AhmanUustalu2016,GambinoKock2013}. A container
is \emph{linear} if it is of the form $S\cdot y$ for a set $S$, i.e.\ its extension is
$\vec X\mapsto S\times X$ in one variable; \emph{finite-support} if its shape set is finite.

\paragraph{Infinitesimal objects and representable tangent bundles.}
We use the representable presentation of tangent structure of \cite{LanfranchiLemay2025},
transported along the dictionary $(\Mod_R,\otimes_R,R)\rightsquigarrow(\Poly,\lhd,y)$.

\begin{definition}[Infinitesimal object]\label{def:inf}
An \emph{infinitesimal object} in $(\Poly,\lhd,y)$ is a polynomial $D$ equipped with an
augmentation: morphisms $z\colon y\to D$ (zero) and $p\colon D\to y$ (projection) with
$p\circ z=\id_y$. The splitting $D=y\oplus M$ defines the \emph{infinitesimal part}
$M=\ker(p)$. Right-$\lhd$-tensoring by $D$ is the endofunctor $T:=(-)\lhd D$, with tangent
projection $A\lhd p\colon A\lhd D\to A$ and zero section $A\lhd z\colon A\to A\lhd D$.
\end{definition}

\begin{definition}[Representable first-order tangent structure]\label{def:tangent}
A \emph{representable first-order tangent structure} is a tangent structure (in the sense of
Cockett--Cruttwell \cite{CockettCruttwell2014}, first order) whose bundle functor is $T=(-)\lhd D$
for an infinitesimal object $D=y\oplus M$. Its infinitesimal part $M$ is called \emph{solid} if
the tangentoid multiplication $\mu\colon M\lhd M\to M$ is an isomorphism.
\end{definition}

The dual-number point $D=y+\varepsilon v$ (read as the infinitesimal object of the base change
$\N\to\N[\varepsilon]/\varepsilon^2$) is such an object: the augmentation is $\varepsilon\mapsto0$,
the zero is $y\hookrightarrow D$, and $(-)\lhd D$ is exactly the dual-number substitution $T$ of
\cite{MacBethTangent}. Its infinitesimal part is the linear container $M=\langle\varepsilon
v\rangle=1\cdot y$, of rank $1$; the trivial structure is $D=y$, $M=0$, $T=\id$. That $T=(-)\lhd D$
is a genuine tangent structure --- not merely an endofunctor --- is the content of
\cite{MacBethTangent}; we take representability of $\partial$ as established there and ask only
whether other infinitesimal objects give rise to further structures.

\paragraph{Why the infinitesimal part is linear.}
For a tangent structure the projection $p\colon T\Rightarrow\id$ is an \emph{additive bundle}: each
fibre carries a commutative-monoid structure and the bundle addition is linear. For $T=(-)\lhd D$
this forces the infinitesimal part $M$ to be a \emph{linear} container $M=S\cdot y$. A non-linear
$M$ --- for instance $\varepsilon y^2$, whose $\varepsilon$-coefficient is $y^2 p'$ rather than a
directional $v\cdot p'$ --- has no vector-space fibre and fails the additive-bundle axiom. So
linearity of $M$ is \emph{forced}, not assumed, and it is this linearity (not any projective-module
theory) that supplies freeness in \S\ref{sec:rigidity}.

\section{The bundle and module tensors coincide on the linear layer}\label{sec:coincide}

The port of the Lanfranchi--Lemay argument has one notorious hazard: the tangent \emph{bundle} is
built from substitution $\lhd$, whereas the \emph{solid} condition $M\otimes M\cong M$ lives in the
``module layer'' where rank multiplies, and these look like two different monoidal products. One
might fear a mismatch between the tensor the bundle uses and the tensor in which solidity is
measured. The following observation dissolves the fear.

\begin{lemma}[Tensor coincidence on the linear layer]\label{lem:coincide}
For linear containers, substitution and the Dirichlet product agree:
\[
  (S\cdot y)\lhd(T\cdot y)\ =\ (S\times T)\cdot y\ =\ (S\cdot y)\otimes_{\mathrm{Dir}}(T\cdot y).
\]
\end{lemma}

\begin{proof}
The extension of $S\cdot y$ is $X\mapsto S\times X$. Composing,
$(S\cdot y)\lhd(T\cdot y)\colon X\mapsto S\times(T\times X)=(S\times T)\times X=(S\times T)\cdot y$.
The Dirichlet product of $\sum_{s}y^1$ and $\sum_t y^1$ is $\sum_{(s,t)}y^{1\times1}=(S\times
T)\cdot y$. Both equal $(S\times T)\cdot y$.
\end{proof}

\begin{remark}
The two tensors agree on monomials ($y^a\lhd y^b=y^{ab}=y^a\otimes_{\mathrm{Dir}}y^b$), and in
particular on the linear layer, which is all the rigidity needs: the infinitesimal part is linear
(\S\ref{sec:prelim}), and $M\lhd M$ is computed there. On polynomials with several monomials the
two diverge --- already $(y+1)\lhd(y+1)=y+2$ while $(y+1)\otimes_{\mathrm{Dir}}(y+1)=y+3$ --- so
Lemma~\ref{lem:coincide} is genuinely a statement about the linear layer, not a global identity.
Neither tensor is the ``doubling'' Dirichlet action $(2y)\otimes_{\mathrm{Dir}}(-)$, which sends
$p\mapsto2p$ and has nothing to do with the derivative.
\end{remark}

\emph{Consequence.} On the linear layer the \emph{bundle} tensor $\lhd$ already multiplies rank, so
the solid condition $M\lhd M\cong M$ reads $|S|^2=|S|$ in the very tensor the bundle uses. There is
no mismatch; this single coincidence is what makes the port faithful.

\section{Rank rigidity}\label{sec:rigidity}

We now classify the finite solid infinitesimal parts. Two structural facts about linear containers,
then the rank computation.

\begin{lemma}[Freeness is automatic over $\N$]\label{lem:free}
Every linear container $S\cdot y$ is the free $\N$-module on the set $S$; two linear containers are
isomorphic iff their shape sets are equinumerous. The cardinal $|S|=:\rank(M)$ is a complete
isomorphism invariant.
\end{lemma}

\begin{proof}
Under coproduct of containers and the $\N$-action $n\cdot(S\cdot y)=(n\times S)\cdot y$, the linear
containers form the free-$\N$-module functor applied to sets: $S\cdot y$ is the $|S|$-fold copower of
$y$. A container isomorphism $S\cdot y\cong T\cdot y$ is, since morphisms between linear containers
are just functions on shapes (positions being singletons), a bijection $S\cong T$. Hence $\rank$ is
a complete invariant.
\end{proof}

This is the analogue of Lanfranchi--Lemay's ``projective $\Rightarrow$ free $\Rightarrow$ rank is a
complete invariant'' (their Lemma~4.16, via Rotman's structure theorem, which \emph{requires} the
base to be a PID). Here it is \emph{automatic}: the linear polynomial functors \emph{are} the free
$\N$-modules by definition, so no PID hypothesis is needed and the rigidity below holds over the
rig $\N$ --- strictly outside the Lanfranchi--Lemay hypotheses.

\begin{lemma}[Solid $\Rightarrow$ idempotent rank]\label{lem:solid}
If $M=S\cdot y$ is solid, i.e.\ $\mu\colon M\lhd M\to M$ is an isomorphism, then $|S|^2=|S|$.
\end{lemma}

\begin{proof}
By Lemma~\ref{lem:coincide}, $M\lhd M=(S\times S)\cdot y$, of rank $|S|^2$ by Lemma~\ref{lem:free}.
An isomorphism $\mu\colon(S\times S)\cdot y\cong S\cdot y$ forces, by Lemma~\ref{lem:free}, a
bijection $S\times S\cong S$, i.e.\ $|S|^2=|S|$.
\end{proof}

\begin{theorem}[Finite rigidity]\label{thm:rigidity}
A finite-support solid infinitesimal part has rank in $\{0,1\}$. Hence there are exactly two:
$M=0$ (rank $0$) and $M=y$ (rank $1$). By \S\ref{sec:prelim} these are the trivial structure and
the AAGM dual-number structure, with combinator $\partial$.
\end{theorem}

\begin{proof}
For a natural number $n=|S|$ we have $n^2=n\iff n(n-1)=0\iff n\in\{0,1\}$, since $\N$ has no zero
divisors. By Lemma~\ref{lem:solid} a finite solid part has $|S|^2=|S|$, so $|S|\in\{0,1\}$; by
Lemma~\ref{lem:free} the rank determines $M$, giving $M\in\{0,y\}$.
\end{proof}

This is the Lanfranchi--Lemay rank computation $R^{n^2}\cong R^n\Rightarrow n^2=n\Rightarrow
n\in\{0,1\}$, ported verbatim with the cardinality rank, and --- per Lemma~\ref{lem:free} --- without
their PID hypothesis.

\paragraph{Machine certification of the arithmetic.}
The finite arithmetic of Theorem~\ref{thm:rigidity} is certified in a Mathlib-free Lean~4
development with no classical axioms. The oracle proves: $n^2=n\iff n\in\{0,1\}$ on $\N$ (in
general, \texttt{omega}-free); that the solid ranks in $0\le n\le12$ are exactly $\{0,1\}$ (by
\texttt{decide}); that these two ranks map through the structure assignment to exactly $\{\id,
\partial\}$, with the dual-number object of arity $2$; that every rank $r\ge2$ (the Whitney sums
of \S\ref{sec:correspondence}) is non-solid; and that the Cantor pairing is injective on a
$12\times12$ window, witnessing $|\N\times\N|=|\N|$ for \S\ref{sec:sharp}. \texttt{\#print axioms}
on the packaged oracle reports only \texttt{propext} (a structural, non-classical axiom); the
finite \texttt{decide} claims are individually axiom-free, with no \texttt{sorryAx},
\texttt{Quot.sound} or \texttt{Classical.choice}. The Lean certifies the \emph{arithmetic collapse}
only --- not the structural correspondence of \S\ref{sec:correspondence}, which remains the open
part.

\section{The finiteness hypothesis is sharp}\label{sec:sharp}

Theorem~\ref{thm:rigidity} is stated for finite support, and the hypothesis cannot be dropped.

\begin{proposition}[Infinite-support obstruction]\label{prop:infinite}
The countable linear container $M=\N\cdot y$ is solid: $M\lhd M=(\N\times\N)\cdot y\cong\N\cdot
y=M$, because $|\N\times\N|=|\N|$. It is a solid infinitesimal part distinct from $0$ and $y$.
Consequently unqualified uniqueness over all of $\Poly$ is false.
\end{proposition}

\begin{proof}
Lemma~\ref{lem:coincide} gives $M\lhd M=(\N\times\N)\cdot y$, and the Cantor pairing supplies a
bijection $\N\times\N\cong\N$ (injectivity checked on a $40\times40$ window; see the Lean oracle of
\S\ref{sec:rigidity} for a certified $12\times12$ window). By Lemma~\ref{lem:free} this is an
isomorphism $\mu\colon M\lhd M\cong M$, and $\rank M=|\N|\notin\{0,1\}$.
\end{proof}

More generally every infinite cardinal $\kappa$ satisfies $\kappa^2=\kappa$ (under choice), so each
yields a solid infinitesimal part. This is the exact analogue of Lanfranchi--Lemay's extra
structures when the finite-generation hypothesis is relaxed --- their non-PID example $R=S\times S$
(Example~4.21) and non-finitely-generated $\mathbb{Z}\ltimes\mathbb{Q}$ (Example~4.20). The moral is
that the honest hypothesis is a \emph{finite-dimensional} infinitesimal part: finitely many tangent
directions, matching Lanfranchi--Lemay's finite-generation (Niefield coexponentiability) condition,
here realised as plain finiteness of the shape set.

\section{From solid parts to tangent structures}\label{sec:correspondence}

Theorem~\ref{thm:rigidity} classifies the finite solid infinitesimal \emph{parts}: exactly
$\{0,y\}$. To pass from parts to tangent \emph{structures} one needs the correspondence
\begin{quote}
$(\star)$\quad representable first-order tangent structure $T=(-)\lhd D$ $\iff$ solid infinitesimal
part $M=\ker(D\to y)$, with the augmented square-zero structure.
\end{quote}
This is Lanfranchi--Lemay's Theorem~3.9 (tangentoid $\iff$ representable tangent structure) composed
with their Theorems~4.8/4.10 (such a tangentoid's infinitesimal part is solid). Over a commutative
ring the correspondence is an equivalence; over the rig $\N$ it splits honestly into a forward
direction we can verify and a converse we cannot yet reprove.

\paragraph{Forward: solid $\Rightarrow$ structure.}
A finite solid part yields a tangent structure in the two surviving cases. $M=0$ gives $T=\id$;
$M=y$ gives the AAGM dual-number structure, which \emph{is} a tangent structure --- the full
Cartesian differential / tangent verification is the content of \cite{MacBethTangent}. Since
Theorem~\ref{thm:rigidity} leaves no third finite solid part, the forward direction never has to
manufacture a third structure.

\paragraph{Converse: structure $\Rightarrow$ solid.}
This is the heart of Lanfranchi--Lemay Theorem~4.8: the \emph{universality of the vertical lift}
forces the tangentoid multiplication to be invertible (solid). Their proof is module-theoretic but
is packaged over a commutative ring; the step that repackages the tangentoid as a square-zero
commutative algebra (their Theorem~4.10) uses ring subtraction, which $\N$ lacks. We therefore
\emph{do not} reprove the converse rig-internally. We state it as the open problem that stands
between Theorem~\ref{thm:rigidity} and a full uniqueness theorem.

\begin{problem}[The solidity converse]\label{prob:converse}
Let $T=(-)\lhd D$ be a representable first-order tangent structure on finite-support polynomial
functors, with infinitesimal part $M=\ker(D\to y)$. Show, internally to $(\Poly,\lhd,y)$ over the
rig $\N$, that the universality of the vertical lift forces $\mu\colon M\lhd M\to M$ to be an
isomorphism (i.e.\ $M$ is solid).
\end{problem}

We offer two pieces of evidence that Problem~\ref{prob:converse} resolves affirmatively.

\emph{Evidence 1 (the non-solid candidates are not primitive structures).} The finite candidate
objects not covered by Theorem~\ref{thm:rigidity} are the higher square-zero objects
\[
  D_r\ =\ y+\varepsilon_1 v_1+\cdots+\varepsilon_r v_r,\qquad \varepsilon_i\varepsilon_j=0,\quad r\ge2,
\]
with infinitesimal part $M=\N^r$, so $M\lhd M=\N^{r^2}$ of rank $r^2\ne r$: \emph{not solid}. But
these $D_r$ are precisely the Whitney sums $T_r=T\times_X\cdots\times_X T$, which occur in the
tangent axioms as \emph{auxiliary} bundles, not as standalone tangent structures. So the non-solid
finite candidates are exactly the ones that are not primitive tangent structures: solidity coincides
with indecomposability, as the correspondence predicts. (The non-solidity of every $r\ge2$ is part
of the certified Lean oracle of \S\ref{sec:rigidity}.)

\emph{Evidence 2 (the module layer is the one the converse uses).} The module layer of polynomial
functors is the free-$\N$-module layer, carrying the same symmetric monoidal structure
(Lemma~\ref{lem:coincide}) and augmentation that Lanfranchi--Lemay's Theorem~4.8 uses. The place
their argument invokes ring subtraction is the $\mathsf{CALG}$ repackaging (Theorem~4.10), not the
solidity derivation (Theorem~4.8), which is module-theoretic throughout. This is why we expect a
rig-internal proof to exist; constructing it is Problem~\ref{prob:converse}.

Granting $(\star_c)$ --- the converse --- we obtain the conditional uniqueness theorem.

\begin{theorem}[Conditional uniqueness]\label{thm:conditional}
Assume the solidity converse of Problem~\ref{prob:converse}. Then among representable first-order
tangent structures $T=(-)\lhd D$ on finite-support polynomial functors there are exactly two: the
trivial structure ($D=y$, $T=\id$) and the AAGM dual-number structure ($D=y+\varepsilon v$,
$T(p)=p(y+\varepsilon v)\bmod\varepsilon^2$, combinator $\partial$).
\end{theorem}

\begin{proof}
By the converse, every such structure has a solid infinitesimal part; by
Theorem~\ref{thm:rigidity} the finite solid parts are exactly $\{0,y\}$; by the forward direction
these are realised by exactly $\{\id,\partial\}$, and by Lemma~\ref{lem:free} the part determines
the object $D=y\oplus M$ up to isomorphism. Hence exactly two structures.
\end{proof}

We state Theorem~\ref{thm:conditional} as explicitly conditional, and we do not claim it
unconditionally: the unconditional half is Theorem~\ref{thm:rigidity}, the classification of solid
parts, which stands on its own.

\section{Comparison with Lanfranchi--Lemay, and significance}\label{sec:compare}

\begin{table}[h]
\centering
\begin{tabular}{p{0.46\textwidth}p{0.46\textwidth}}
\toprule
Lanfranchi--Lemay (affine schemes) & This note (polynomial functors) \\
\midrule
base a commutative ring $R$; rigidity needs $R$ a \textbf{PID} &
  base the rig $\N$ (no subtraction); rigidity needs \textbf{no} PID \\
projective $\Rightarrow$ free $=$ Rotman structure theorem (uses PID) &
  linear functor \emph{is} free $\N$-module (Lemma~\ref{lem:free}) --- automatic \\
rank in $R$-modules, $R^{n^2}\cong R^n$ &
  cardinality rank, $(S\times S)\cong S$ (Lemmas~\ref{lem:coincide},~\ref{lem:free}) \\
solid $M\otimes_R M\cong M\Rightarrow n^2=n\Rightarrow n\in\{0,1\}$ &
  solid $M\lhd M\cong M\Rightarrow|S|^2=|S|\Rightarrow|S|\in\{0,1\}$ \\
two structures: trivial $+$ Kähler/dual-number &
  two parts: $0+y$, realising $\id+\partial$ \\
extras at non-PID / non-f.g.\ (Ex.~4.20--4.21) &
  extra at \textbf{infinite support} $\N\cdot y$ (\S\ref{sec:sharp}) \\
\bottomrule
\end{tabular}
\caption{The port is faithful (same rank-idempotent mechanism) and sharper on one axis: freeness is
automatic over $\N$, so the rigidity lives over a rig rather than a PID. The tensor coincidence
(Lemma~\ref{lem:coincide}) is the linchpin specific to polynomials.}
\label{tab:compare}
\end{table}

The classification is a clean deliverable for two audiences. For the tangent-category community
\cite{CockettCruttwell2014,LanfranchiLemay2025} it exhibits a rig-based model in which the
representable-tangent rigidity holds for a strictly weaker reason than over a PID, and isolates the
single structural step (Problem~\ref{prob:converse}) that a rig-internal uniqueness theorem still
needs. For the \texttt{Poly} and containers community \cite{AAGM2003,AhmanUustalu2016} it upgrades
``the container derivative is \emph{a} tangent structure'' \cite{MacBethTangent} toward ``it is the
\emph{only} nontrivial one,'' quantifying exactly how much of that upgrade is proved (the solid-part
classification) and how much is reduced to a named lemma (the vertical-lift converse). The shape of
the target --- \emph{we classified the tangent structures on finite-support polynomial functors} ---
is the clean statement; this note delivers its rigidity half.

\section{Conclusion}\label{sec:conclusion}

We have proved that a finite-support solid infinitesimal part of $(\Poly,\lhd,y)$ has rank in
$\{0,1\}$, hence is either $0$ or $y$, realising the trivial structure and the AAGM container
derivative; the arithmetic is Lean-certified with no classical axioms. The one technical
observation that makes the port of Lanfranchi--Lemay's affine-scheme rigidity faithful is that on
the linear layer the bundle tensor $\lhd$ and the module tensor coincide
(Lemma~\ref{lem:coincide}), so solidity is measured in the tensor the bundle already uses; and the
port is cleaner than its template because over the rig $\N$ freeness is automatic, with no
principal-ideal hypothesis. The finiteness is sharp: $\N\cdot y$ is a third solid part, so
unqualified uniqueness over $\Poly$ fails.

What remains is a single, precisely located step. Full uniqueness
(Theorem~\ref{thm:conditional}) follows once one knows that every representable first-order tangent
structure has a solid infinitesimal part --- the rig-internal analogue of Lanfranchi--Lemay's
vertical-lift universality (Problem~\ref{prob:converse}). We have given two pieces of evidence that
it holds: the non-solid finite candidates are exactly the Whitney sums, which are auxiliary rather
than primitive; and the derivation that needs the converse is module-theoretic, not ring-theoretic,
so the absence of subtraction in $\N$ should not obstruct it. Three routes to close the gap suggest
themselves: a direct Poly-internal proof of the vertical-lift universality; the operadic machinery
of Rosický tangent categories of algebras over an operad, the species-adjacent setting closest to
polynomial functors; and making ``finite support $\iff$ coexponentiable in $(\Poly,\lhd)$'' precise,
so the finiteness hypothesis of Theorem~\ref{thm:rigidity} becomes intrinsic rather than imposed.
Higher Weil algebras $\N[\varepsilon]/\varepsilon^k$ ($k\ge3$) and iterated tangents are first-order
out of scope, as they are for the derivative, for Lanfranchi--Lemay, and for the fibrational theory
\cite{CCGZ2024}.

\begin{thebibliography}{9}

\bibitem{AAGM2003}
M.~Abbott, T.~Altenkirch, N.~Ghani, C.~McBride,
\emph{Derivatives of containers},
Typed Lambda Calculi and Applications (TLCA 2003), LNCS 2701, Springer, 2003, 16--30.

\bibitem{AhmanUustalu2016}
D.~Ahman, T.~Uustalu,
\emph{Directed containers as categories},
Proc.\ MSFP 2016, EPTCS 207, 2016, 89--98.

\bibitem{BCS2009}
R.~F.~Blute, J.~R.~B.~Cockett, R.~A.~G.~Seely,
\emph{Cartesian differential categories},
Theory and Applications of Categories \textbf{22} (2009), no.~23, 622--672.

\bibitem{CCGZ2024}
M.~Capucci, G.~Cruttwell, N.~Ghani, F.~Zanasi,
\emph{A fibrational theory of first order differential structures},
preprint, arXiv:2409.05763 (2024).

\bibitem{CockettCruttwell2014}
J.~R.~B.~Cockett, G.~S.~H.~Cruttwell,
\emph{Differential structure, tangent structure, and SDG},
Applied Categorical Structures \textbf{22} (2014), no.~2, 331--417.

\bibitem{GambinoKock2013}
N.~Gambino, J.~Kock,
\emph{Polynomial functors and polynomial monads},
Math.\ Proc.\ Cambridge Philos.\ Soc.\ \textbf{154} (2013), no.~1, 153--192.

\bibitem{LanfranchiLemay2025}
M.~Lanfranchi, J.-S.~P.~Lemay,
\emph{Representable tangent structures for affine schemes},
preprint, arXiv:2505.09080 (2025).

\bibitem{MacBethTangent}
MacBeth,
\emph{The container derivative is a tangent structure: a fibrational positioning}
(Kodamai internal note, 2026).

\end{thebibliography}

\end{document}
